% Page budget: 1.35 pages | Owns: gantry coordinates, first-principles equations, scheduling, LPV-LFR realisation, and physical parameterisation.
% WRITING GUIDE
% Purpose: Define the physical coordinates and derive the exact continuous-time LPV-LFR baseline used by every later section.
% Primary reference for Section II-B and its appendix material:
% docs/Thesis-documentation/LPV_LFR_Stepwise_Derivation/main.tex
% (the version shared with Jasper, including the G-Delta interconnection figure).
% Follow its signal definitions, direct mass-equation construction, block partitioning,
% loop-closing verification, and allocation of detailed algebra to the appendix.
% Supporting checks/details only: LPV_LFR_internal_loop_Well_posedness,
% LPV_Mass_invertible, LPV_LFR_Justification_Drenth, LPV_LFR_Derivation,
% and docs/fp-model-structure.md. Where their presentation differs, follow
% LPV_LFR_Stepwise_Derivation for Section II-B.
% Sources: README "Source map per subsection", rows II-A, II-B, II-C and II Frozen-LTI baseline; mandatory papers and the reference gate as in the README.
% Research-plan anchor: Preserve Aspect 1's explicit position dependence; update the original discretisation wording to the final method.
% Current decisions: Continuous-time physics, endogenous scheduler Y, fixed-step RK4, and identifiable parameter combinations.
% Choices to justify: Coordinate frame, Y as scheduler, continuous-time formulation, RK4 step and substeps, full six-channel LFR realization, and reduced physical parameterization. Give the physical or numerical reason for each choice, not only its definition.
% Cautions: The final pipeline uses the full six-channel scheduling loop, not the experimental six-to-four SVD reduction. Resolve the supervisor comment about dual-gantry terminology; do not copy unverified bibliography entries.
% Planned figures: Physical coordinate schematic and compact interconnection between G and Delta(Y).
% Section is finished when: The LFR collapses to the original model, well-posedness assumptions are stated, and notation matches Section 03.
%
% LPV-LFR DERIVATION ROADMAP
% 1. Define q=[X,Theta,Y]^T, the force transformation, measured outputs, assumptions, and the admissible range of Y.
% 2. State M(Y)qdd+Cqd+Kq=f and expose the structural dependence M(Y)=M0+Y M1+Y^2 M2.
% 3. Construct the LFR directly from the polynomial mass equation:
%       a1=Y qdd, a2=Y a1, M0 qdd + M1 a1 + M2 a2 = f_net,
%       z=[qdd; a1], w=[a1; a2] (stepwise latent names; qdd is not renamed to a).
% 4. Collect w=Delta(Y)z with Delta(Y)=Y I6 and give the constant block-level G realization with all signal dimensions.
% 5. Eliminate w and z in a short exactness check to recover the original equation of motion.
% 6. State the well-posedness condition over the complete scheduling range and relate it to invertibility of M(Y); place the longer proof outside the main text.
% 7. Only after the conceptual construction, mention that the implementation
%    evaluates the resulting rational expression in Horner form. Put
%    M(Y)^(-1)=[N0+Y N1+Y^2 N2]/d(Y), the explicit Ni entries, and the
%    determinant coefficients in the appendix or technical supplement.
% 8. Keep this derivation continuous-time; Section III-A owns the RK4 scheme, Section V only its step size.
%
% DERIVATION WARNING
% Do not start from v=M(Y)^(-1)f_net. Follow LPV_LFR_Stepwise_Derivation:
% construct the repeated scheduling loop from qdd, a1=Y qdd, a2=Y a1 with
% z=[qdd; a1], w=[a1; a2], then verify it by closing the loop. Older inverse-first notes
% may be used only as supporting algebra. Verify terminology against the
% primary literature by Drenth and Hoekstra.
% MUST ESTABLISH (II-B after eq:delta, agreed 2026-10-06):
% 1. G is a constant LTI system in Drenth Eq. (6) form; its first block row is the state equation.
% 2. Closing the algebraic loop gives back eq:eom exactly.
% 3. Well posed iff M(Y) nonsingular, true for every Y; hence closed-form M(Y)^{-1}=N(Y)/d(Y)
%    (N_i, d(Y) and Horner evaluation in the appendix). The reduced-parameter bound stays in II-C.
% MUST ESTABLISH (II-C, agreed 2026-10-07):
% 1. The fourteen raw parameters enter M0,M1,M2,C,K only through theta_base (eq:phi); this work's own.
% 2. Argument in the main text, ordered toward the conclusion: input-output behaviour fixes the matrices;
%    twelve entries, three equal m_h, ten independent; so exactly ten combinations are structurally
%    identifiable; only then named theta_base (eq:phi). Model-only (no data), no new symbols, no
%    Proposition/proof environment (agreed 2026-10-07).
% 3. Four lost directions in one sentence; estimates reported in theta_base.
% 4. Positive theta_base does not give M(Y)>0; that holds iff eq:lfr_admissibility (proof in app:lfr-wellposed).
% 5. One sentence: training keeps eq:lfr_admissibility at every iterate. The m_Delta map (eq:mdelta_map)
%    and its justification are stated in Section III-C (agreed 2026-10-07), not here.
% Compact form agreed 2026-10-07: no separate definition paragraph, no null-direction display.
% Moved to Section V as bullets: rho value, combination-error metric with the m_Delta normalisation, raw vs
% reduced coordinates per arm, practical identifiability on the records. Appendix app:params keeps the table only.
% Main text: assumptions, polynomial dependence, direct LFR construction, compact G, exactness, well-posedness statement, and one implementation sentence.
% Appendix/supplement: expanded M_i and N_i matrices, determinant polynomial, longer invertibility/well-posedness proof, and symbolic verification.
%
% DRAFT NOTES (cycle 12, 2026-10-09)
% MUST ESTABLISH (II-A, PROPOSED, inferred from the README ownership and source-map rows):
% 1. Coordinates q=[X,Theta,Y]^T and the P map between actuator and generalized frames, with the reason (Garcia Sec. 2.2).
% 2. EOM under Garcia's small-yaw and no-Coriolis assumptions; Y makes the baseline quasi-LPV (Toth Def. 7.2).
% 3. Adapted, not adopted: Y kept (Garcia Eq. 11 drops it), Coulomb terms of Garcia Eq. 6 omitted (todo D-022 vs D-204).
% MUST ESTABLISH (II-D, PROPOSED, from README ownership and source-map row II Frozen-LTI):
% 1. M_LPV and the frozen M_LTI named together, input u_act, output q_act, nominal parameters, Y_op value in Section V.
% Corrections to the agreed blocks: II-B item 1 cites the continuous-time form of drenth2025thesis Eqs. (2.1), (2.2),
%   since drenth2025lpvlfr Eq. (6) is discrete time and names its LTI part M (README source map II-B).
%   II-B item 3: written with adj M(Y) and det M(Y), since d already names the payload offset and N(Y)/d(Y)
%   would give d two meanings; "true for every Y" holds at physical parameter values (app:lfr-wellposed).
%   II-C item 2: hong2020global is not cited beside ovchinnikov2021computing (README source map II-C);
%   gautier1990minimum dropped (no local PDF, unverified), see todo.
% Left to other sections: RK4 step and substeps (Section III-A owns the scheme, Section V the step size);
%   the m_Delta training map (Section III-C); the operating range of the records (Section V).
% Existing todos resolved by this draft: dynamics outside the baseline (now the opening's hand-over to the
%   augmentation); coordinate reason and "how to connect" (Garcia Sec. 2.2 sentence and the Y paragraph);
%   "source, or state earlier" (Toth Def. 7.2); "look at this section" and "left here" (style remarks);
%   geometric dimensions and Telica identification (values are Garcia Table I; real data are out of the
%   main results); SVD reduction (excluded: header Caution, Meeting-audit candidate-thesis-impact.md).
% Stale lines elsewhere (not edited): 03_augmentation.tex line 361 cites eq:wellposed, now defined here;
%   01_introduction.tex Aspect 1 todo ("LPV versus frozen LTI exists only as a baseline check") and
%   06_results.tex lines 18 and 23 now have M_LPV, M_LTI to name; 99_appendix remarks list moved into II-B.
\section{Dual-Drive Gantry System and Physics Baseline}\label{sec:baseline}
\ifdraft
\begin{itemize}
\item Opening: the augmentation needs an exact, well-posed baseline in identifiable parameters; names II-A to II-D; delivers the first contribution of Section~I (Introduction; style profile)
\end{itemize}
\fi

The augmentation of Section~I starts from a first-principles baseline of the
gantry, which must be exact, well posed for every payload position and
parameterised by quantities the data determine.
To this end, we state the coordinates and equations of motion
(Section~\ref{sec:sys}) and realise them as a continuous-time LPV-LFR
(Section~\ref{sec:lfr}).
We then derive the identifiable parameter combinations
(Section~\ref{sec:params}) and define the two baseline models of the
comparison (Section~\ref{sec:baseline_models}).
The realisation is the first contribution of Section~I: a compact
continuous-time LPV-LFR of the dual-drive gantry baseline with rational
dependence on the payload position.
\todo{Missing: contribution wording. This section also derives the ten identifiable combinations and the exact admissibility condition \eqref{eq:lfr_admissibility}, which the first contribution does not name. Candidate: ``Derive a compact continuous-time LPV-LFR realisation of the dual-drive gantry baseline with rational dependence on payload position, its exact well-posedness condition and its identifiable parameter combinations'' (README contribution map; D-190, D-229).}
\todo{Missing: source of the supervisor comment on dual-gantry terminology (header Caution; no decision entry or meeting-audit item found). Candidate: Garc\'ia-Herreros et al.'s term ``dual-drive gantry'' throughout, as in this section's title; Section~I still writes ``dual-gantry'' (\cite{garcia2013model}, title and Sec.~2.2).}

\subsection{System and equations of motion}\label{sec:sys}
\ifdraft
\begin{itemize}
\item Generalized coordinates $(X,\Theta,Y)$ avoid a closed kinematic chain and expose the load-distribution coupling; $P$ maps the actuator frame of the data to them (Garc\'ia Sec.~2.2; adopted)
\item $L_b$ fixed, so the coordinates do not depend on estimated parameters (code; display)
\item Small yaw and no Coriolis terms give $M(Y)\ddot q+C\dot q+Kq=P\uact$; $Y$ is a state, so the baseline is quasi-LPV (Garc\'ia Sec.~2.3; T\'oth Def.~7.2)
\item Adapted: $Y$ kept, unlike Garc\'ia's control model (11); Coulomb terms of Garc\'ia (6) omitted (?) (D-022 vs D-204)
\item Fourteen raw parameters $\theta$, nominal values from Garc\'ia Table~I in Appendix~\ref{app:params}
\end{itemize}
\fi

The baseline needs coordinates in which the equations of motion are compact
and the dependence on the payload position is explicit.
We adopt the lumped-parameter model of the dual-drive H-type gantry of
Garc\'ia-Herreros et al.~\cite[Sec.~2.2]{garcia2013model}
(Fig.~\ref{fig:gantry_system}).
Two parallel linear actuators drive a cross-arm through flexible joints.
A third actuator moves the payload along the cross-arm.
The model uses the generalized coordinates $q=[X\;\Theta\;Y]^\top$, the
translation and yaw of the cross-arm and the payload position, rather than
the actuator positions $\qact=[X_1\;X_2\;Y]^\top$.
These coordinates avoid a closed kinematic chain and expose the coupling
caused by the non-uniform load distribution over the two drives.

\begin{figure}[b]
  \centering
  \def\gantryabsorber{0}% baseline only; the absorber version is in Section III
  \input{figures/gantry_system_with_absorber}
  \caption{Lumped-parameter model of the dual-drive gantry, adapted from
  \cite{garcia2013model}. Red labels denote actuator coordinates and blue
  labels generalized coordinates.}
  \label{fig:gantry_system}
\end{figure}

The data hold the actuator positions and the actuator forces
$\uact=[F_{X1}\;F_{X2}\;F_Y]^\top$, which map to the generalized frame as
\begin{equation}
\begin{aligned}
  q &= P^{-\top}\qact,\qquad u = P\uact,\\
  P &=
  \begin{bmatrix}
    1 & 1 & 0\\
    L_b/2 & -L_b/2 & 0\\
    0 & 0 & 1
  \end{bmatrix},
\end{aligned}
  \label{eq:Ptransform}
\end{equation}
where $\qact,\uact,q\in\R^3$, $u\in\R^3$ is the generalized force, and
$L_b$ is the cross-arm length.
$L_b$ is fixed at its nominal value, so $P$ and $q$ do not depend on the
estimated parameters.

Moving the payload changes the mass distribution, so the equations of
motion carry a position-dependent inertia.
Following Garc\'ia-Herreros et al.~\cite[Sec.~2.3]{garcia2013model}, we set
$\cos\Theta\approx1$ and $\sin\Theta\approx0$, since the yaw stays within
tens of microradians.
We also neglect the Coriolis and centripetal terms.
The equations of motion are then
\begin{equation}
  M(Y)\ddot{q}(t) + C\dot{q}(t) + K q(t) = P\uact(t),
  \label{eq:eom}
\end{equation}
{\footnotesize
\setlength{\arraycolsep}{2pt}
\begin{equation}
M(Y)=
\begin{bmatrix}
m_1+m_2+m_b+m_h & \tfrac{L_b}{2}(m_1-m_2)-m_hY & 0\\
\tfrac{L_b}{2}(m_1-m_2)-m_hY &
\substack{J_b+J_h+\tfrac{L_b^2}{4}(m_1+m_2)\\{}+m_h(d^2+Y^2)} & -m_hd\\
0 & -m_hd & m_h
\end{bmatrix},
\label{eq:mass_matrix}
\end{equation}
\begin{equation}
\begin{aligned}
C&=
\begin{bmatrix}
c_{g1}+c_{g2} & \tfrac{L_b}{2}(c_{g1}-c_{g2}) & 0\\
\tfrac{L_b}{2}(c_{g1}-c_{g2}) &
c_{b1}+c_{b2}+\tfrac{L_b^2}{4}(c_{g1}+c_{g2}) & 0\\
0 & 0 & c_y
\end{bmatrix},\\[1ex]
K&=
\begin{bmatrix}
0&0&0\\
0&k_{b1}+k_{b2}&0\\
0&0&0
\end{bmatrix},
\end{aligned}
\label{eq:damping_stiffness_matrices}
\end{equation}
}
where $m_1$, $m_2$, $m_b$ and $m_h$ are the masses of the two actuators,
the cross-arm and the payload, $J_b$ and $J_h$ the yaw inertias of the
cross-arm and the payload, $d$ the offset of the payload from the cross-arm,
$c_{g1}$, $c_{g2}$ and $c_y$ the actuator damping coefficients, and
$c_{b1}$, $c_{b2}$, $k_{b1}$ and $k_{b2}$ the damping and stiffness of the
flexible joints, and $\theta\in\R^{14}$ collects these fourteen parameters
(nominal values from~\cite[Table~I]{garcia2013model} in
Appendix~\ref{app:params}).
Since $Y$ is a state, \eqref{eq:eom} is a quasi-LPV model scheduled by
$Y$~\cite[Def.~7.2]{toth2010modeling}.

The baseline adapts the model of Garc\'ia-Herreros et al.\ in two respects.
It keeps $Y$ as a coordinate, which their control model removes to leave
$X$ and $\Theta$~\cite[Eq.~(11)]{garcia2013model}.
It omits the Coulomb friction they model at the three
actuators~\cite[Eq.~(6)]{garcia2013model}.
\todo{Missing: reason for omitting the Coulomb terms. D-204 keeps friction in the simulated system only, so that the augmentation must capture it; D-022 requires the baseline to be the first-principles model as derived, which includes these terms. Candidate: the D-204 reason, stated once Section~V has introduced the friction of the simulated system (D-204, D-224).}

\subsection{Self-scheduled LPV-LFR representation}\label{sec:lfr}
\ifdraft
\begin{itemize}
\item The state equation of \eqref{eq:eom} needs $M(Y)^{-1}$, rational in $Y$; an LPV-LFR moves this dependence into $\Dsch(Y)$; the scheduling map is known, unlike Drenth (Drenth thesis Sec.~2.1; Drenth et al.\ Sec.~3.2)
\item $M(Y)=M_0+YM_1+Y^2M_2$; $Y^2$ realised by applying $Y$ twice, as in Drenth's rational embedding (Drenth thesis Sec.~2.1.1; adapted)
\item Latent variables $z,w$ and Delta-block $\Dsch(Y)=YI_6$; six channels, not claimed minimal (this work; header Caution)
\item Continuous time, unlike Drenth et al., so $\theta$ keeps the structure of $M(Y)$, $C$, $K$ (D-018, D-020)
\item LTI part $G$ in the continuous-time form of Drenth thesis Eqs.~(2.1), (2.2); the dependence on $Y$ is static (adapted; research-plan anchor)
\item Closing the algebraic loop gives back \eqref{eq:eom} exactly (this work)
\item Well posed iff $M(Y)$ nonsingular, exact unlike Drenth et al.\ Thm.~6; true for every $Y$ at physical values (Drenth thesis Sec.~2.1; this work)
\item The code solves the loop in closed form; the LFR stays the form of Section~III's baseline block (D-020; code)
\end{itemize}
\fi

The state equation of \eqref{eq:eom} needs $M(Y)^{-1}$, whose entries are
rational in $Y$.
An LPV-LFR isolates such a dependence: a constant LTI part $G$ is closed by
a scheduling block $\Dsch(Y)$ (Fig.~\ref{fig:lfr})~\cite[Sec.~2.1]{drenth2025thesis}.
Here the scheduling variable is the state $Y$.
The scheduling map is therefore known and, unlike in the self-scheduled
models of Drenth et al.~\cite[Sec.~3.2]{drenth2025lpvlfr}, not estimated.

\begin{figure}[t]
  \centering
  \input{figures/lfr_loop}
  \caption{Self-scheduled baseline interconnection. The Delta-block closes
  the latent variables of $G$ through $w=\Dsch(Y)z$, with
  $\Dsch(Y)=YI_6$ and $Y$ a component of the state.}
  \label{fig:lfr}
\end{figure}

We construct the realisation directly from the mass equation.
The inertia matrix \eqref{eq:mass_matrix} is quadratic in $Y$,
\begin{equation}
  M(Y)=M_0+YM_1+Y^2M_2,
  \label{eq:Msplit}
\end{equation}
where $M_0=M(0)$ and $M_1,M_2\in\R^{3\times3}$ hold the $m_h$ terms
(Appendix~\ref{app:lfr-details}).

Following the rational embedding of Drenth~\cite[Sec.~2.1.1]{drenth2025thesis},
we realise $Y^2$ by applying $Y$ twice.
A separate scheduling variable for $Y^2$ would admit values that cannot
occur.
With $a_1=Y\ddot q$ and $a_2=Ya_1$, \eqref{eq:eom} becomes the $Y$-free
relation
\begin{equation}
  M_0\ddot q+M_1a_1+M_2a_2=u-C\dot q-Kq,
  \label{eq:lfr_constant}
\end{equation}
where $a_1,a_2\in\R^3$.
The two multiplications by $Y$ form the scheduling loop
\begin{equation}
  w=\begin{bmatrix}a_1\\a_2\end{bmatrix}
   =\begin{bmatrix}YI_3&0\\0&YI_3\end{bmatrix}
    \begin{bmatrix}\ddot q\\a_1\end{bmatrix}
   =\Dsch(Y)z,
  \label{eq:delta}
\end{equation}
where $z,w\in\R^6$ are the latent variables and $\Dsch(Y)=YI_6$.
This six-channel realisation is not claimed to be minimal.

The realisation stays in continuous time, unlike the discrete-time LPV-LFR
of Drenth et al.~\cite[Sec.~3.1]{drenth2025lpvlfr}.
In continuous time, $\theta$ keeps the structure of $M(Y)$, $C$ and $K$,
which a discretisation before identification would not preserve.
Solving \eqref{eq:lfr_constant} for $\ddot q$ with $M_0^{-1}$ gives the LTI
part in the continuous-time form of Drenth~\cite[Eqs.~(2.1), (2.2)]{drenth2025thesis},
{\small
\begin{equation}
\begin{aligned}
  \begin{bmatrix}\dot{\tilde x}\\ z\\ q\end{bmatrix}
  &=\underbrace{\begin{bmatrix}
    A_x&B_w&B_u\\ C_z&D_{zw}&D_{zu}\\ C_y&0&0
  \end{bmatrix}}_{G}
  \begin{bmatrix}\tilde x\\ w\\ u\end{bmatrix},\\
  D_{zw}&=\begin{bmatrix}-M_0^{-1}M_1&-M_0^{-1}M_2\\ I_3&0\end{bmatrix},
\end{aligned}
  \label{eq:lfr_G}
\end{equation}
}
where $\tilde x=\col(q,\dot q)\in\R^6$ is the state, $G\in\R^{15\times15}$,
and the remaining blocks are given in Appendix~\ref{app:lfr-details}.
The scheduling dependence is static: only $\Dsch$ depends on $Y$, and not
on its derivatives.

Closing the loop with \eqref{eq:delta} gives back $a_1=Y\ddot q$ and
$a_2=Y^2\ddot q$, so \eqref{eq:lfr_constant} returns \eqref{eq:eom}.
The LFR is therefore an exact realisation of the baseline.

The LFR is well posed if its loop has a unique solution for every $Y$, that
is, if $I_6-D_{zw}\Dsch(Y)$ is nonsingular~\cite[Sec.~2.1]{drenth2025thesis}.
Drenth et al.\ give a sufficient condition that bounds the scheduling set and
the spectral radius of $D_{zw}$~\cite[Thm.~6]{drenth2025lpvlfr}.

For the gantry realisation, we derive the exact condition instead.
A Schur complement gives (Appendix~\ref{app:lfr-details})
\begin{equation}
  \det\bigl(I_6-D_{zw}\Dsch(Y)\bigr)=\frac{\det M(Y)}{\det M_0},
  \label{eq:wellposed}
\end{equation}
so the LFR is well posed exactly where $M(Y)$ is nonsingular.
For positive masses and yaw inertias, $M(Y)\succ0$ for every
$Y\in\R$ (Appendix~\ref{app:lfr-wellposed}).
\todo{Missing: operating range of $Y$ (mechanical stroke, simulated range, range of the records). Candidate: no range is needed for the baseline, since it is well posed for every $Y$; Section~V states the range of the records (DATA-DESIGN.md Sec.~7).}

The implementation solves the loop in closed form,
$\ddot q=M(Y)^{-1}(u-C\dot q-Kq)$, with $\operatorname{adj}M(Y)$ and
$\det M(Y)$ evaluated as quadratics in $Y$ by Horner's scheme.
The LFR remains the form of the baseline block in Section~III, which
evaluates the first block row of $G$ with $w=\Dsch(Y)z$.

\subsection{Physical parameters and identifiable combinations}\label{sec:params}
\ifdraft
\begin{itemize}
\item The input-output behaviour fixes $M_0,M_1,M_2,C,K$ and conversely, so the identifiable part of $\theta$ is what the matrices contain (this work)
\item Twelve nonzero entries, three equal to $m_h$, leave ten independent functions: exactly ten combinations are structurally identifiable (this work; definition Ovchinnikov et al.\ Def.~7)
\item Written physically as $\theta_{\mathrm{base}}$; four lost directions; estimates reported in $\theta_{\mathrm{base}}$ (this work)
\item Positive $\theta_{\mathrm{base}}$ does not give $M(Y)\succ0$; that holds iff \eqref{eq:lfr_admissibility} (this work; proof Appendix~\ref{app:lfr-wellposed})
\item Training keeps \eqref{eq:lfr_admissibility} at every iterate; the map and its justification are in Section~\ref{sec:aug_closed_loop}
\end{itemize}
\fi

Estimating the baseline from data needs parameters that the data determine.
We derive which combinations of the fourteen entries of $\theta$ qualify.
With all three positions measured and all three generalized forces actuated,
every value of $(Y,q,\dot q,\ddot q)$ is attained.
Two parameter vectors with equal input-output behaviour therefore give equal
$M(Y)\ddot q+C\dot q+Kq$ at every such value.
Hence they give equal $M_0$, $M_1$, $M_2$, $C$ and $K$.
Conversely, these matrices fix the behaviour, so the identifiable part of
$\theta$ is what they contain.

Up to symmetry, \eqref{eq:mass_matrix} and \eqref{eq:damping_stiffness_matrices}
have twelve nonzero entries.
Three of them, $M_{0,33}$, $-M_{1,12}$ and $M_{2,22}$, equal $m_h$, which
leaves ten distinct functions of $\theta$.
Their Jacobian with respect to
$(m_h,d,m_1,m_b,J_b,c_{g1},c_{g2},c_{b1},c_y,k_{b1})$ has a determinant of
magnitude $m_hL_b^2/2\neq0$, so the ten functions are independent.
Exactly ten combinations of $\theta$ are therefore structurally
identifiable~\cite[Def.~7]{ovchinnikov2021computing}.
Four directions of $\theta$ are lost: $k_{b1}-k_{b2}$, $c_{b1}-c_{b2}$,
$J_b-J_h$, and a transfer of mass from the actuators to the cross-arm with
compensating $J_b+J_h$.
We write the ten combinations in physically readable form as
\begin{equation}
  \theta_{\mathrm{base}}=[k_{b,\Sigma},c_{g1},c_{g2},c_y,c_{b,\Sigma},
  m_h,m_\Sigma,m_\Delta,J_{\mathrm{eff}},d]^\top\in\R^{10},
  \label{eq:phi}
\end{equation}
where $k_{b,\Sigma}=k_{b1}+k_{b2}$, $c_{b,\Sigma}=c_{b1}+c_{b2}$,
$m_\Sigma=m_1+m_2+m_b$, $m_\Delta=m_1-m_2$ and
$J_{\mathrm{eff}}=J_b+J_h+\tfrac{L_b^2}{4}(m_1+m_2)$.
Estimates are therefore reported in $\theta_{\mathrm{base}}$.
\todo{Missing: (a) verified source for the analogy with the base inertial parameters of robot identification. Candidate: Gautier and Khalil 1990 (refs.bib key gautier1990minimum; no local PDF, unverified), else leave out. (b) Notation: the index of $m_\Delta$ and $\xi_\Delta$ reuses $\Delta$, which names the scheduling block; both symbols are used in Sections~III and~IV. Candidate: $m_{\mathrm{d}}$ (Dirk decides).}

Positive combinations do not imply a positive definite inertia.
For $m_h,m_\Sigma,J_{\mathrm{eff}}>0$, $M(Y)\succ0$ for every $Y$ if and
only if
\begin{equation}
  m_\Sigma J_{\mathrm{eff}}>\tfrac{L_b^2}{4}\,m_\Delta^2
  \label{eq:lfr_admissibility}
\end{equation}
(Appendix~\ref{app:lfr-wellposed}).
Section~\ref{sec:aug_closed_loop} trains $\theta_{\mathrm{base}}$ in
coordinates that keep \eqref{eq:lfr_admissibility}, so every iterate gives a
well-posed LFR.

\subsection{Baseline models}\label{sec:baseline_models}
\ifdraft
\begin{itemize}
\item Aspect~1 compares the position-dependent with a position-independent baseline, so both are needed as models (Introduction Aspect~1; THESIS-RESULTS step~1)
\item $\mathcal M_{\mathrm{LPV}}$ and frozen $\mathcal M_{\mathrm{LTI}}$: input $\uact$, output $\qact$, nominal $\theta$; $Y_{\mathrm{op}}$ value in Section~V (README ownership; D-242; notation (?))
\item Compared without training as a demonstration of the size of the position-dependent error (THESIS-RESULTS step~1)
\end{itemize}
\fi

Aspect~1 of Section~I asks whether the position-dependent form predicts
better than a position-independent one, so the comparison needs both forms
as models.
Both take the actuator forces and return the actuator positions, with
$\theta$ at the nominal values of Appendix~\ref{app:params}:
\begin{equation}
\begin{aligned}
  \mathcal M_{\mathrm{LPV}}&:\;
  M(Y)\ddot q+C\dot q+Kq=P\uact,\quad \qact=P^\top q,\\
  \mathcal M_{\mathrm{LTI}}&:\;
  M(Y_{\mathrm{op}})\ddot q+C\dot q+Kq=P\uact,\quad \qact=P^\top q,
\end{aligned}
\label{eq:baseline_models}
\end{equation}
where $Y_{\mathrm{op}}\in\R$ is a fixed payload position, whose value
Section~\ref{sec:setup} gives.
Section~\ref{sec:results} compares the two without training, as a demonstration of the
size of the position-dependent error.
\todo{Missing: (a) model names. Candidate: the calligraphic $\mathcal M_{\mathrm{LPV}}$, $\mathcal M_{\mathrm{LTI}}$, since plain $M$ clashes with $M(Y)$ (README open conflicts). (b) The value of $Y_{\mathrm{op}}$ is not yet in Section~V. Candidate: $Y_{\mathrm{op}}=0$, mid-stroke, since $Y$ is measured from the cross-arm centre (THESIS-RESULTS step~1; \cite[Sec.~2.2]{garcia2013model}). (c) No runner evaluates $\mathcal M_{\mathrm{LTI}}$ yet; the code provides it as the frozen branch of the baseline block (README source map, II Frozen-LTI).}
